\section{Signale und Systeme}
\label{sec:signals_and_systems}
In diesem Abschnitt werden einige für den weiteren Verlauf wichtige Begriffe definiert, insbesondere \textit{Signal} und \textit{System},
sowie die sogenannten Linearen Zeit-Invarianten Systeme.
\paragraph{Signale}
Ein Signal beschreibt die Veränderung eines Parameters im Verhältnis zu einem anderen Parameter.
Signale übertragen Informationen über den Zustand oder das Verhalten eines physischen Systems. Mathematisch sind Signale
Funktionen einer oder mehrerer unabhängiger Variablen. Die unabhängige Variable von Audiosignalen repräsentiert Zeit oder Frequenz \cite[Kap. 5]{smith1997}.
Die unabhängige Variable kann stetig
oder diskret sein. Stetige Signale werden oft als \textit{analoge} Signale bezeichnet \cite[Kap. 2]{oppenheim1999}. In Computergestützten Anwendungen werden 
ausschließlich diskrete Signale verarbeitet.
Der Definitionsbereich diskreter Signale sind die natürlichen Zahlen. Diskrete Signale sind also mathematisch betrachtet Folgen. 
Signale $S:\mathbb{N}\mapsto\mathbb{Z}$ heißen \textit{digitale Signale} \cite[Kap. 2]{oppenheim1999}.
Stetige Signale werden mit der unabhängigen Variable in runden Klammern notiert:
$x(t)$ für ein stetiges Signal $x$ in Abhängigkeit von Zeit $t$.
Diskrete Signale werden mit der unabhängigen Variablen in eckigen Klammern notiert: $x[n]$ für ein diskretes Signal $x$,
wobei $n$ den Index des betrachteten Funktionswerts bezeichnet \cite[Seite 9]{oppenheim1999}.
\paragraph{Der Einheitsimpuls}
Ein besonderes diskretes Signal ist die Einheits-Sample-Folge.
Sie ist definiert als die Folge
\begin{equation}
    \delta[n]=
    \begin{cases}
        0,\,n\ne0\\
        1,\,n=0
    \end{cases}
\end{equation}
Jedes diskrete Signal lässt sich als eine Summe skalierter und verzögerter Einheitsimpulsen darstellen \cite[Kap. 2, S. 11]{oppenheim1999}.
Im weiteren Verlauf wird die Einheits-Sample-Folge etwas vereinfacht als ``Einheitsimpuls'' bezeichnet.
\paragraph{Systeme}
Ein System produziert als Antwort auf ein Eingangssignal ein Ausgangssignal \cite[Kap. 5]{smith1997}.
Formal definiert \cite[Seite 16]{oppenheim1999} ein diskretes System als eine Transformation oder einen Operator, welche(r) ein Eingangssignal $x[n]$
auf ein Ausgangssignal $y[n]$ abbildet. Die Funktion $T:x[n]\mapsto y[n]$ heißt \textit{Transferfunktion} des Systems.
\paragraph{Lineare Zeit-Invariante Systeme (LTI-Systeme)}
Besonders signifikant für die Signalverarbeitung sind Lineare Zeit-Invariante Systeme. Damit ein System \textit{linear}
heißt, muss es die beiden Eigenschaften
\begin{equation}
    T(x_1[n]+x_2[n])=T(x_1[n])+T(x_2[n])=y_1[n]+y_2[n]
\end{equation}
und 
\begin{equation}
    T(a\cdot x[n])=a\cdot T(x[n])=a\cdot y[n]
\end{equation}
erfüllen.
Daraus resultiert die \textit{Superposition} von Eingangssignalen:
\begin{equation}
    x[n]=\sum_{k}a_kx_k[n]\Rightarrow y[n]=\sum_{k}a_ky_k[n]
\end{equation}
Dabei beizeichnet $x[n]$ das Eingangssignal, $y[n]$ das Ausgangssignal, und $a_k$ eine Folge von Konstanten \cite[Kap. 2]{oppenheim1999}

Ein System heißt \textit{Zeit-Invariant}, wenn für jedes $n_0$ und jedes Eingangssignal $x$ mit den Werten
$x'[n]=x[n-n_0]$ das entsprechende Ausgangssignal $y$ mit den Werten $y'[n]=y[n-n_0]$ erzeugt wird \cite[Kap. 2]{oppenheim1999}.

Systeme mit beiden Eigenschaften heißen \textit{lineare Zeit-Invariante Systeme (LTI-Systeme)}.
Die Auswirkungen von LTI-Systemen auf ein beliebiges Eingangssignal können vollständig durch ihre Impulsantwort (engl.: \textit{Impulse Response}, kurz \textit{IR},
siehe \ref{sec:convolution}) ausgedrückt werden.
Die Impulsantwort eines Systems ist selbst ein diskretes Signal und wird mit $h[n]$ notiert \cite[Kap. 2]{oppenheim1999}.
